\vspace{1em}
\section{Related Work}
\label{section:relatedwork}


\noindent{\bf Model serving systems:} Systems such as Clipper~\cite{clipper}, TensorFlow-Serving~\cite{tensorflowserving}, Clockwork~\cite{clockwork} and BatchMaker~\cite{batchmakereurosys} study various placement, caching and batching strategies for model serving. However, these systems fail to address the challenges of auto-regressive transformer inference. More recently, systems such as Orca~\cite{orca}, vLLM~\cite{vllmsosp}, FlexGen~\cite{flexgen}, FasterTransformers~\cite{fastertransformer}, LightSeq~\cite{lightseq}, and TurboTransformers~\cite{turbotransformers}  propose domain-specific optimizations for transformer inference. FlexGen~\cite{flexgen} optimizes LLM inference for throughput in resource-constrained offline scenarios i.e., it is not suitable for online serving. FastServe~\cite{fastserve} proposed a preemptive scheduling framework for LLM inference to minimize the job completion times. We present a detailed comparison with Orca and vLLM as they represent the state-of-the-art in LLM inference. 


Another approach that has emerged recently is to disaggregate the prefill and decode phases on different replicas as proposed in SplitWise, DistServe and TetriInfer ~\cite{patel2023splitwise,distserve2024,tetriinfer}. These solutions can entirely eliminate the interference between prefills and decodes. However, disaggregation requires migrating the KV cache of \textit{each} request upon the completion of its prefill phase which could be challenging in the absence of high-bandwidth interconnects between different replicas. In addition, this approach also under-utilizes the GPU memory capacity of the prefill replicas i.e., only the decode replicas are responsible for storing the KV cache. On the positive side, disaggregated approaches can execute prefills with maximum efficiency (and therefore yield better TTFT) unlike chunked prefills that are somewhat slower than full prefills. We leave a quantitative comparison between \sysname and disaggregation-based solutions for future work.

Recently, Sheng et al. \cite{sheng2023fairness} proposed modification to \ils algorithm to ensure fairness among clients in a multi-tenant environment. FastServe~\cite{fastserve} uses a preemption based scheduling mechanism to mitigate head-of-the-line blocking. Such algorithmic optimizations are complimentary to our approach and can benefit from lower prefill-decode interference enabled by \sysname. Another recent system, APIServe \cite{abhyankar2024apiserve} adopted chunked prefills from \sarathi to utilize wasted compute in decode batches for ahead-of-time prefill recomputation for multi-turn API serving.




\noindent{\bf Improving GPU utilization for transformers:} Recent works have proposed various optimizations to improve the hardware utilization for transformers. FasterTransformer uses model-specific GPU kernel implementations. CocoNet~\cite{coconetasplos} and ~\cite{googleoverlap} aim to overlap compute with communication to improve GPU utilization: these techniques are specially useful while using a high degree of tensor-parallel for distributed models where communication time can dominate compute. Further, the cost of computing self-attention grows quadratically with sequence length and hence can become significant for long contexts. \cite{self:attn:on2:memory, flashattention, flashattention2} have proposed various techniques to minimize the memory bottlenecks of self-attention with careful tiling and work partitioning. In addition, various parallelization strategies have been explore to optimize model placement. These techniques are orthogonal to \sysname.  


\noindent{\bf Model optimizations:} A significant body of work around model innovations has attempted to address the shortcomings of  transformer-based language models or to take the next leap forward in model architectures, beyond transformers. For example, multi-query attention~\cite{multiqueryattention} shares the same keys and values across all the attention heads to reduce the size of the KV-cache, allowing to fit a larger batch size on the GPUs. Several recent works have also shown that the model sizes can be compressed significantly using quantization~\cite{smoothquant,gptq,qlora,llmint8}. Mixture-of-expert models are aimed primarily at reducing the number of model parameters that get activated in an iteration~\cite{largescalemoe,moeatc23,moedeployment}. More recently, retentive networks have been proposed as a successor to transformers~\cite{retnet}. In contrast, we focus on addressing the performance issues of popular transformer models from a GPU's perspective. 
